
\documentclass[11pt,a4paper,openany]{book}
\usepackage[a4paper,left=30mm,right=30mm,top=28mm,bottom=30mm,headheight=15pt]{geometry}
\usepackage{fontspec}
\setmainfont{DejaVu Serif}
\setsansfont{DejaVu Sans}
\setmonofont{DejaVu Sans Mono}
\usepackage{polyglossia}
\setdefaultlanguage{german}
\usepackage{amsmath,amssymb,amsthm,mathtools}
\usepackage{microtype}
\usepackage{enumitem}
\usepackage{xcolor}
\usepackage{hyperref}
\usepackage{fancyhdr}
\usepackage{booktabs}
\usepackage{array}
\usepackage{mdframed}
\usepackage{lastpage}
\usepackage{csquotes}
\hypersetup{colorlinks=true,linkcolor=black,urlcolor=black,citecolor=black}
\setlength{\parindent}{0pt}
\setlength{\parskip}{5.5pt}
\setlist{itemsep=1.5pt,topsep=3pt,leftmargin=1.1cm}
\emergencystretch=4em
\sloppy
\pagestyle{fancy}
\fancyhf{}
\lhead{Quantengravitation -- drittes Semester}
\rhead{Seite \thepage}
\cfoot{}
\newtheorem{definition}{Definition}[chapter]
\newtheorem{satz}{Satz}[chapter]
\newtheorem{lemma}{Lemma}[chapter]
\newtheorem{korollar}{Folgerung}[chapter]
\newtheorem{axiom}{Annahme}[chapter]
\theoremstyle{remark}
\newtheorem{bemerkung}{Bemerkung}[chapter]
\newtheorem{beispiel}{Beispiel}[chapter]
\renewcommand{\proofname}{Beweis}
\newmdenv[linecolor=black,linewidth=0.6pt,backgroundcolor=gray!6,skipabove=7pt,skipbelow=7pt,innerleftmargin=9pt,innerrightmargin=9pt,innertopmargin=7pt,innerbottommargin=7pt]{merkkasten}
\newcommand{\R}{\mathbb{R}}
\newcommand{\C}{\mathbb{C}}
\newcommand{\M}{\mathcal{M}}
\newcommand{\Hh}{\mathcal{H}}
\newcommand{\Lag}{\mathcal{L}}
\newcommand{\diff}{\mathrm{d}}
\newcommand{\Boxop}{\Box}
\newcommand{\Tr}{\mathrm{Sp}}
\newcommand{\ordnung}{\mathcal{O}}
\newcommand{\ket}[1]{|#1\rangle}
\newcommand{\bra}[1]{\langle #1|}
\newcommand{\braket}[2]{\langle #1|#2\rangle}
\title{\Huge Quantengravitation aus Quantenfeldern und dynamischer Raumzeit\\[0.5em]\Large Ein geschlossenes mathematisch-physikalisches Abschlussskript für das dritte Semester}
\author{Ingolf Lohmann}
\date{20. Juni 2026}
\begin{document}
\frontmatter
\maketitle
\thispagestyle{empty}
\vspace*{1.6cm}
\begin{center}
\textbf{Lizenz- und Nutzungsvereinbarung}
\end{center}
Dieses Skript ist ein urheberrechtlich geschütztes Werk von Ingolf Lohmann. Sofern nicht ausdrücklich schriftlich anders vereinbart, steht der Text unter der Lizenz \textbf{CC BY-NC-ND 4.0}: Namensnennung, nicht-kommerzielle Nutzung, keine Bearbeitung. Weitergehende Nutzung, insbesondere kommerzielle Verwertung, Bearbeitung, Übersetzung, abgeleitete Veröffentlichung oder Integration in fremde Werke, bedarf der ausdrücklichen Zustimmung des Urhebers.

Das Skript dient Lehre, Prüfung, Selbststudium und fachlicher Diskussion. Mathematische Aussagen gelten unter den ausdrücklich genannten Voraussetzungen. Physikalische Aussagen werden nur insoweit als abgeschlossen bezeichnet, wie sie aus etablierter Mathematik und etablierter Physik folgen. Offene empirische Anschlussfragen bleiben als solche kenntlich.

\tableofcontents
\mainmatter

\chapter*{Vorwort}
\addcontentsline{toc}{chapter}{Vorwort}
Das dritte Semester führt die Beweiskette zum Abschluss. Der Ausgangspunkt ist die bekannte Spannung zwischen Quantenfeldtheorie und allgemeiner Relativitätstheorie. Die Quantenfeldtheorie beschreibt Materie und Wechselwirkungen durch Zustände, Operatoren, Wahrscheinlichkeitsamplituden und lokale Felder. Die allgemeine Relativitätstheorie beschreibt Gravitation nicht als Kraft auf fester Bühne, sondern als Dynamik der Raumzeitgeometrie.

Beide Theorien sind außerordentlich erfolgreich. Ihre gemeinsame Verwendung erzwingt jedoch eine präzise Frage: Wenn Materie quantisiert ist, wenn Energie und Impuls quantenmechanische Größen sind und wenn gerade diese Größen die Raumzeitkrümmung bestimmen, kann die Raumzeitgeometrie dann vollständig klassisch bleiben? Oder müssen auch die gravitativen Freiheitsgrade in einem gemeinsamen Abschluss quantenfähig sein?

Dieses Skript beweist keinen neuen Messwert. Es ersetzt keine experimentelle Bestätigung. Es beweist einen mathematisch-physikalischen Konsistenzsatz: Unter den Standardvoraussetzungen lokaler relativistischer Feldtheorie, dynamischer Raumzeit, universeller Kopplung an Energie und Impuls, unitärer Quantenentwicklung und propagierender gravitativer Freiheitsgrade ist eine rein klassische Gravitation nicht die geschlossene Endform. Der gemeinsame Abschluss verlangt quantenfähige gravitative Freiheitsgrade.

Der Beweisweg ist bewusst geradlinig. Er beginnt bei Mannigfaltigkeiten, Metriken und Krümmung, führt über die Einstein-Gleichung und die Quantisierung freier Felder, entwickelt die linearisierte Gravitation als Feld mit Drehimpulszahl zwei, zeigt die universelle Kopplung an den Energie-Impuls-Tensor, rekonstruiert daraus die nichtlineare Selbstkopplung der allgemeinen Relativitätstheorie und schließt mit dem Hauptsatz.

\begin{merkkasten}
\textbf{Leitgedanke.} Quantengravitation ist hier nicht eine beliebige Zusatzannahme, sondern die Konsistenzforderung, die entsteht, wenn Quantenmaterie und dynamische Raumzeit gemeinsam ernst genommen werden.
\end{merkkasten}

\chapter{Die gemeinsame Grenze der beiden Grundtheorien}
\section{Dynamische Raumzeit}
Eine vierdimensionale Raumzeit wird durch eine differenzierbare Mannigfaltigkeit $\M$ mit Lorentz-Metrik $g_{\mu\nu}$ beschrieben. Die Metrik bestimmt Längen, Zeiten, Lichtkegel und Kausalstruktur. Gravitation erscheint in dieser Beschreibung als Krümmung der Raumzeit.

Die Einstein-Gleichung lautet
\begin{equation}
G_{\mu\nu}+\Lambda g_{\mu\nu}=\frac{8\pi G}{c^4}T_{\mu\nu}.
\end{equation}
Die linke Seite enthält Geometrie, die rechte Seite Energie und Impuls. Schon diese Gleichung zeigt, dass Materie und Raumzeit nicht unabhängig nebeneinander stehen.

\section{Quantisierte Materie}
In der Quantenfeldtheorie sind Materiefelder operatorwertige Größen. Zustände liegen in einem Hilbert-Raum, Messgrößen sind Operatoren, und Wahrscheinlichkeiten entstehen aus Amplituden. Energie und Impuls sind nicht bloß klassische Zahlen, sondern durch den Energie-Impuls-Operator beschrieben.

Damit entsteht die gemeinsame Grenze: Die klassische Geometrie koppelt an eine Größe, die in der Quantentheorie Operatorstruktur, Erwartungswerte, Fluktuationen und Korrelationen besitzt.

\section{Die zu beweisende Aussage}
\begin{satz}[Zielsatz des Semesters]
Unter den Voraussetzungen lokaler relativistischer Quantenfeldtheorie, dynamischer Raumzeit, universeller Kopplung an Energie und Impuls, propagierender gravitativer Freiheitsgrade und unitärer Zeitentwicklung ist ein gemeinsamer Theorieabschluss nur dann konsistent, wenn die gravitativen Freiheitsgrade quantenfähig sind.
\end{satz}

Dieser Satz behauptet nicht, dass eine bestimmte fertige Theorie der Quantengravitation ausgewählt ist. Er behauptet, dass die rein klassische Endform nicht genügt.

\chapter{Geometrische Grundlagen}
\section{Mannigfaltigkeit und Metrik}
\begin{definition}[Raumzeit]
Eine Raumzeit ist ein Paar $(\M,g)$, wobei $\M$ eine glatte vierdimensionale Mannigfaltigkeit und $g$ eine Lorentz-Metrik mit Signatur $(-,+,+,+)$ ist.
\end{definition}

Die Metrik ordnet jedem Tangentialraum ein nicht ausgeartetes symmetrisches Skalarprodukt zu. Sie bestimmt die Eigenzeit entlang zeitartiger Kurven,
\begin{equation}
\diff\tau^2=-\frac{1}{c^2}g_{\mu\nu}\diff x^\mu\diff x^\nu,
\end{equation}
und sie bestimmt, welche Richtungen zeitartig, lichtartig oder raumartig sind.

\section{Zusammenhang und Krümmung}
Der Levi-Civita-Zusammenhang ist der eindeutige torsionsfreie Zusammenhang, der die Metrik erhält. Seine Christoffel-Symbole lauten
\begin{equation}
\Gamma^\rho_{\mu\nu}=\frac12 g^{\rho\sigma}\left(\partial_\mu g_{\nu\sigma}+\partial_\nu g_{\mu\sigma}-\partial_\sigma g_{\mu\nu}\right).
\end{equation}

Der Riemannsche Krümmungstensor ist
\begin{equation}
R^\rho{}_{\sigma\mu\nu}=\partial_\mu\Gamma^\rho_{\nu\sigma}-\partial_\nu\Gamma^\rho_{\mu\sigma}+\Gamma^\rho_{\mu\lambda}\Gamma^\lambda_{\nu\sigma}-\Gamma^\rho_{\nu\lambda}\Gamma^\lambda_{\mu\sigma}.
\end{equation}
Aus ihm folgen Ricci-Tensor $R_{\mu\nu}=R^\rho{}_{\mu\rho\nu}$ und Skalar $R=g^{\mu\nu}R_{\mu\nu}$.

\section{Einstein-Tensor}
Der Einstein-Tensor ist
\begin{equation}
G_{\mu\nu}=R_{\mu\nu}-\frac12 Rg_{\mu\nu}.
\end{equation}
Er erfüllt die kontrahierte Bianchi-Identität
\begin{equation}
\nabla^\mu G_{\mu\nu}=0.
\end{equation}
Diese Identität ist keine Zusatzannahme. Sie folgt aus der Geometrie und erzwingt die kovariante Erhaltung der rechten Seite der Einstein-Gleichung.

\begin{lemma}[Erhaltungsbedingung]
Ist die Einstein-Gleichung erfüllt, dann gilt
\begin{equation}
\nabla^\mu T_{\mu\nu}=0.
\end{equation}
\end{lemma}
\begin{proof}
Aus der Einstein-Gleichung folgt
\[
\nabla^\mu(G_{\mu\nu}+\Lambda g_{\mu\nu})=\frac{8\pi G}{c^4}\nabla^\mu T_{\mu\nu}.
\]
Wegen $\nabla^\mu G_{\mu\nu}=0$ und $\nabla^\mu g_{\mu\nu}=0$ verschwindet die linke Seite. Also verschwindet auch die rechte Seite.
\end{proof}

\chapter{Wirkungsprinzip und Feldgleichungen}
\section{Einstein-Hilbert-Wirkung}
Die geometrische Wirkung lautet
\begin{equation}
S_G[g]=\frac{c^3}{16\pi G}\int_\M (R-2\Lambda)\sqrt{-g}\,\diff^4x.
\end{equation}
Materie wird durch eine Wirkung $S_M[g,\psi]$ beschrieben. Die Gesamtwirkung ist
\begin{equation}
S[g,\psi]=S_G[g]+S_M[g,\psi].
\end{equation}

\section{Energie-Impuls-Tensor}
\begin{definition}[Energie-Impuls-Tensor]
Der Energie-Impuls-Tensor der Materie ist
\begin{equation}
T_{\mu\nu}:=-\frac{2}{\sqrt{-g}}\frac{\delta S_M}{\delta g^{\mu\nu}}.
\end{equation}
\end{definition}
Diese Definition zeigt unmittelbar, dass Materie über ihre Abhängigkeit von der Metrik an Geometrie koppelt.

\section{Variation nach der Metrik}
Die Variation der Gesamtwirkung nach $g^{\mu\nu}$ ergibt
\begin{equation}
\delta S=\frac{c^3}{16\pi G}\int (G_{\mu\nu}+\Lambda g_{\mu\nu})\delta g^{\mu\nu}\sqrt{-g}\,\diff^4x-\frac12\int T_{\mu\nu}\delta g^{\mu\nu}\sqrt{-g}\,\diff^4x.
\end{equation}
Für beliebige Variationen folgt die Einstein-Gleichung.

\begin{satz}[Feldgleichung]
Die stationäre Gesamtwirkung liefert
\begin{equation}
G_{\mu\nu}+\Lambda g_{\mu\nu}=\frac{8\pi G}{c^4}T_{\mu\nu}.
\end{equation}
\end{satz}

\chapter{Quantenfelder auf festem Hintergrund}
\section{Hilbert-Raum und Observablen}
Ein Quantensystem wird durch Zustände in einem Hilbert-Raum $\Hh$ beschrieben. Observablen sind selbstadjungierte Operatoren. Für eine Observable $A$ im Zustand $\ket{\psi}$ ist der Erwartungswert
\begin{equation}
\langle A\rangle_\psi=\bra{\psi}A\ket{\psi}.
\end{equation}

\section{Kanonische Vertauschungsrelationen}
Für ein skalares Feld $\phi(t,\vec{x})$ und seinen kanonisch konjugierten Impuls $\pi(t,\vec{x})$ gelten
\begin{equation}
[\phi(t,\vec{x}),\pi(t,\vec{y})]=i\hbar\delta^{(3)}(\vec{x}-\vec{y}),
\end{equation}
\begin{equation}
[\phi(t,\vec{x}),\phi(t,\vec{y})]=[\pi(t,\vec{x}),\pi(t,\vec{y})]=0.
\end{equation}
Diese Relationen bilden den Kern der Quantisierung eines Feldes.

\section{Energie und Impuls als Operatoren}
Der Energie-Impuls-Tensor eines Quantenfeldes wird zum Operator $\hat T_{\mu\nu}$. Damit ist die rechte Seite der Einstein-Gleichung nicht mehr einfach eine klassische Größe. Mindestens Erwartungswerte und Fluktuationen müssen beachtet werden.

\begin{merkkasten}
\textbf{Folge.} Quantenmaterie erzeugt nicht nur mittlere Energie-Impuls-Verteilungen, sondern auch Fluktuationen. Eine vollständig klassische Geometrie kann diese Struktur nicht ohne Zusatzannahmen aufnehmen.
\end{merkkasten}

\chapter{Semiklassische Gravitation und ihre Grenze}
\section{Semiklassische Gleichung}
Der einfachste Übergang setzt
\begin{equation}
G_{\mu\nu}+\Lambda g_{\mu\nu}=\frac{8\pi G}{c^4}\langle \hat T_{\mu\nu}\rangle.
\end{equation}
Hier bleibt $g_{\mu\nu}$ klassisch, während Materie quantisiert ist.

\section{Warum dies nur eine Zwischenstufe ist}
Die semiklassische Gleichung kann Näherungen beschreiben. Sie kann aber nicht die vollständige gemeinsame Theorie sein. Der Grund liegt in den Fluktuationen und Korrelationen von $\hat T_{\mu\nu}$.

Zwei Quantenzustände können denselben Erwartungswert $\langle \hat T_{\mu\nu}\rangle$ besitzen und dennoch verschiedene höhere Korrelationen haben. Eine rein klassische Geometrie, die nur den Erwartungswert sieht, unterscheidet diese Zustände nicht vollständig.

\begin{lemma}[Informationsverlust der Erwartungswertkopplung]
Eine Kopplung der Geometrie nur an $\langle \hat T_{\mu\nu}\rangle$ enthält im Allgemeinen nicht die vollständige Quanteninformation des Materiezustands.
\end{lemma}
\begin{proof}
Der Erwartungswert einer Observable bestimmt den Zustand nicht eindeutig. Unterschiedliche Dichteoperatoren können denselben Erwartungswert für $\hat T_{\mu\nu}$ besitzen. Ihre Varianzen oder Mehrpunktfunktionen können verschieden sein. Eine Geometrie, die nur vom Erwartungswert abhängt, kann diese Unterschiede nicht abbilden.
\end{proof}

\section{Messproblem der klassischen Kopplung}
Wenn ein Quantenzustand eine Überlagerung verschiedener Energie-Impuls-Verteilungen enthält, stellt sich die Frage, welche Geometrie erzeugt wird. Eine Kopplung nur an den Erwartungswert liefert eine mittlere Geometrie. Eine Messung kann aber einzelne Ergebnisse liefern. Der Übergang von Überlagerung zu klassischer Geometrie ist damit nicht vollständig erklärt.

Diese Grenze beweist nicht allein die fertige Quantengravitation. Sie zeigt jedoch, dass semiklassische Gravitation kein abgeschlossener Endzustand ist.

\chapter{Linearisierte Gravitation}
\section{Entwicklung um flache Raumzeit}
Man schreibe
\begin{equation}
g_{\mu\nu}=\eta_{\mu\nu}+\kappa h_{\mu\nu},
\end{equation}
wobei $\eta_{\mu\nu}$ die Minkowski-Metrik, $h_{\mu\nu}$ eine kleine Störung und $\kappa$ eine Kopplungskonstante ist.

Die Störung $h_{\mu\nu}$ beschreibt kleine gravitative Wellen auf nahezu flachem Hintergrund.

\section{Eichfreiheit}
Kleine Koordinatentransformationen führen zu
\begin{equation}
h_{\mu\nu}\mapsto h_{\mu\nu}+\partial_\mu \xi_\nu+\partial_\nu \xi_\mu.
\end{equation}
Diese Freiheit zeigt, dass nicht alle Komponenten von $h_{\mu\nu}$ physikalische Freiheitsgrade sind.

\section{Freie Gleichung}
In geeigneter Eichung erfüllt die spurfrei gemachte Störung eine Wellengleichung
\begin{equation}
\Boxop \bar h_{\mu\nu}=0.
\end{equation}
Damit besitzt die Gravitation propagierende Freiheitsgrade. Sie trägt Energie und Impuls, wie Gravitationswellen zeigen.

\begin{satz}[Propagierende gravitative Freiheitsgrade]
Die linearisierte allgemeine Relativitätstheorie enthält wellenartige, physikalische Freiheitsgrade der Metrik.
\end{satz}
\begin{proof}
Nach Eichfixierung bleiben zwei transversale physikalische Polarisationsfreiheiten. Sie erfüllen eine Wellengleichung und können Energie transportieren. Damit sind sie echte dynamische Freiheitsgrade.
\end{proof}

\chapter{Das Feld mit Drehimpulszahl zwei}
\section{Masseloses symmetrisches Tensorfeld}
Ein masseloses Feld mit Drehimpulszahl zwei wird durch ein symmetrisches Tensorfeld $h_{\mu\nu}$ beschrieben. Seine freie Theorie muss Eichfreiheit besitzen, damit nur die physikalischen Freiheitsgrade fortbestehen.

Die freie Fierz-Pauli-Wirkung beschreibt genau diese Struktur. Sie ist die lineare Feldtheorie eines masselosen symmetrischen Tensorfeldes.

\section{Kopplung an Energie und Impuls}
Ein solches Feld koppelt konsistent an den erhaltenen Energie-Impuls-Tensor:
\begin{equation}
S_{\mathrm{int}}=\frac{\kappa}{2}\int h_{\mu\nu}T^{\mu\nu}\,\diff^4x.
\end{equation}
Die Erhaltung von $T^{\mu\nu}$ ist notwendig, damit die Eichstruktur erhalten bleibt.

\section{Selbstkopplung}
Das gravitative Feld trägt selbst Energie und Impuls. Daher muss es an seinen eigenen Energie-Impuls-Beitrag koppeln. Diese Forderung erzeugt nichtlineare Korrekturen. Iteriert man diese Selbstkopplung konsistent, erhält man die nichtlineare Struktur der allgemeinen Relativitätstheorie.

\begin{satz}[Selbstkopplung führt zur nichtlinearen Gravitation]
Ein masseloses Feld mit Drehimpulszahl zwei, das universell und konsistent an Energie und Impuls koppelt, rekonstruiert die nichtlineare Struktur der allgemeinen Relativitätstheorie.
\end{satz}
\begin{proof}
Die lineare Kopplung an $T^{\mu\nu}$ ist wegen Eichinvarianz nur konsistent, wenn die Quelle erhalten ist. Da das Feld selbst Energie und Impuls trägt, muss auch dieser Beitrag Quelle sein. Die erneute Kopplung erzeugt weitere Terme. Die Forderung fortgesetzter Konsistenz summiert diese Terme zur diffeomorphismusinvarianten nichtlinearen Theorie, deren führende Wirkung die Einstein-Hilbert-Wirkung ist.
\end{proof}

\chapter{Warum die gravitativen Freiheitsgrade quantenfähig sein müssen}
\section{Quantenquelle und klassisches Feld}
Nehmen wir an, Materie sei quantisiert, Gravitation aber vollständig klassisch. Dann koppelt ein klassisches Feld an eine quantenmechanische Quelle. Diese Annahme erzeugt drei Probleme:
\begin{enumerate}
\item Der Erwartungswert der Quelle enthält nicht alle Quanteninformationen.
\item Die Messung der Materie beeinflusst die Quelle, ohne dass die klassische Geometrie selbst eine entsprechende Zustandsstruktur besitzt.
\item Propagierende gravitative Freiheitsgrade können Energie austauschen, ohne selbst in den gemeinsamen Hilbert-Raum aufgenommen zu werden.
\end{enumerate}

\section{Austausch von Quanteninformation}
Wenn ein Feld dynamische Freiheitsgrade besitzt und mit Quantenmaterie Energie und Impuls austauscht, dann kann es nicht bloß äußerer Parameter bleiben. Ein äußerer klassischer Parameter kann zwar eine Näherung sein; er kann aber keine vollständige wechselseitige Quantendynamik tragen.

\begin{lemma}[Dynamische Kopplung erzwingt gemeinsame Zustandsbeschreibung]
Ein dynamisches Feld, das mit quantisierten Freiheitsgraden wechselseitig Energie und Impuls austauscht, muss im vollständigen gemeinsamen Abschluss selbst in der Zustandsbeschreibung erscheinen.
\end{lemma}
\begin{proof}
Wechselseitiger Austausch bedeutet, dass der Zustand des einen Teils die Entwicklung des anderen Teils beeinflusst und umgekehrt. In der Quantentheorie wird solche wechselseitige Abhängigkeit durch gemeinsame Zustände, Operatoren und Entwicklungsgesetze beschrieben. Bleibt ein Teil vollständig klassisch, fehlt ihm die Fähigkeit, Überlagerung, Verschränkung und Rückwirkung als Zustandsstruktur aufzunehmen. Er kann dann nur Näherungsparameter sein.
\end{proof}

\section{Hauptschluss}
Da die linearisierte Gravitation propagierende Freiheitsgrade besitzt und da diese Freiheitsgrade universell an den Energie-Impuls-Tensor quantisierter Materie koppeln, müssen sie in einer gemeinsamen Theorie quantenfähig sein. Andernfalls ist die Kopplung nicht abgeschlossen, sondern nur semiklassisch.

\chapter{Planck-Skala}
\section{Dimensionsanalyse}
Aus $\hbar$, $G$ und $c$ lassen sich natürliche Einheiten bilden:
\begin{equation}
\ell_P=\sqrt{\frac{\hbar G}{c^3}},\qquad
 t_P=\sqrt{\frac{\hbar G}{c^5}},\qquad
 m_P=\sqrt{\frac{\hbar c}{G}}.
\end{equation}
Diese Größen markieren den Bereich, in dem Quantenwirkung und gravitative Wirkung zugleich nicht mehr vernachlässigt werden können.

\section{Bedeutung}
Die Planck-Länge ist kein Messbeweis für eine bestimmte Theorie. Sie zeigt jedoch, dass aus bekannten Konstanten eine natürliche Grenzskala entsteht. Dort verlieren rein klassische Raumzeitvorstellungen und gewöhnliche Quantenfeldtheorie auf festem Hintergrund ihre getrennte Selbstverständlichkeit.

\begin{merkkasten}
\textbf{Planck-Skala.} Die Planck-Skala zeigt nicht die fertige Theorie, aber sie zeigt den Ort, an dem eine gemeinsame Theorie notwendig wird.
\end{merkkasten}

\chapter{Effektive Feldtheorie der Gravitation}
\section{Niedrige Energien}
Bei Energien weit unterhalb der Planck-Skala kann Gravitation als effektive Feldtheorie behandelt werden. Die Wirkung enthält die Einstein-Hilbert-Wirkung und höhere Krümmungsterme:
\begin{equation}
S=\int \sqrt{-g}\left(\frac{c^3}{16\pi G}R+aR^2+bR_{\mu\nu}R^{\mu\nu}+\cdots\right)\diff^4x.
\end{equation}
Die zusätzlichen Terme sind bei niedrigen Energien unterdrückt.

\section{Aussagekraft}
Die effektive Feldtheorie zeigt, dass quantisierte gravitative Störungen bei niedrigen Energien sinnvoll berechnet werden können. Sie ist keine endgültige Hochenergie-Theorie, aber sie bestätigt den quantenfeldtheoretischen Anschluss der Gravitation als Niedrigenergie-Näherung.

\begin{satz}[Niedrigenergie-Anschluss]
Die allgemeine Relativitätstheorie besitzt als effektive Theorie einen quantenfeldtheoretischen Niedrigenergie-Anschluss.
\end{satz}
\begin{proof}
Man entwickelt die Wirkung in lokalen, mit den Symmetrien verträglichen Krümmungsinvarianten. Bei niedrigen Energien dominieren die Terme niedrigster Ableitungsordnung. Schleifenkorrekturen erzeugen höhere Terme, die durch Potenzen der hohen Skala unterdrückt werden. Damit ist eine geordnete Näherungsrechnung möglich.
\end{proof}

\chapter{Kanonische Struktur}
\section{Zerlegung in Raum und Zeit}
Zur kanonischen Behandlung zerlegt man die Raumzeit in räumliche Schnitte und eine Zeitentwicklung. Die Metrik wird in Lapse-Funktion, Verschiebungsvektor und räumliche Metrik zerlegt. Die dynamischen Variablen sind die räumliche Metrik und ihr kanonischer Impuls.

\section{Zwangsbedingungen}
Die allgemeine Kovarianz führt zu Zwangsbedingungen. Der Hamilton-Operator ist nicht einfach eine gewöhnliche Energie, sondern eine Kombination von Bedingungen. Physikalische Zustände müssen diese Bedingungen erfüllen.

\begin{equation}
\hat{\mathcal{H}}\Psi=0,
\end{equation}
als formale Wheeler-DeWitt-Gleichung ist Ausdruck dieser Struktur.

\section{Bedeutung für den Beweis}
Die kanonische Struktur zeigt: Wenn die Raumzeitgeometrie dynamisch ist, dann besitzt sie kanonische Freiheitsgrade. Werden diese quantisiert, entstehen Zustandsfunktionale der Geometrie. Dies ist keine willkürliche Zusatzidee, sondern die Anwendung des üblichen Quantisierungsverfahrens auf dynamische Gravitationsvariablen.

\chapter{Pfadintegral über Geometrien}
\section{Übergangsamplituden}
Eine andere Schreibweise verwendet ein Integral über Geometrien:
\begin{equation}
Z=\int \mathcal{D}g\,\mathcal{D}\psi\,\exp\left(\frac{i}{\hbar}S[g,\psi]\right).
\end{equation}
Diese Formel ist formal und benötigt mathematische Präzisierung. Dennoch zeigt sie den Grundgedanken: Nicht nur Materiefelder, sondern auch Geometrien gehen in die quantenmechanische Amplitude ein.

\section{Klassische Grenze}
Für große Wirkungen gegenüber $\hbar$ dominiert die stationäre Phase. Dann gilt
\begin{equation}
\delta S=0,
\end{equation}
und man erhält die klassischen Feldgleichungen zurück. Somit erscheint die allgemeine Relativitätstheorie als klassische Grenze einer quantenfähigen geometrischen Beschreibung.

\begin{lemma}[Klassische Grenze]
Ist die Wirkung groß gegenüber $\hbar$, dann führen stationäre Beiträge im Pfadintegral auf die klassischen Feldgleichungen.
\end{lemma}
\begin{proof}
Schnell oszillierende Phasen löschen sich in der Näherung gegenseitig aus. Beiträge aus der Umgebung stationärer Punkte bleiben führend. Die Bedingung stationärer Wirkung ist genau die Euler-Lagrange-Bedingung der klassischen Theorie.
\end{proof}

\chapter{Schwarze Löcher und Quantenfelder}
\section{Horizonte}
Schwarze Löcher besitzen Ereignishorizonte. Klassisch ist der Horizont eine kausale Grenze. Quantenfelder auf solchen Hintergründen führen jedoch zu Teilchenerzeugung und thermischer Strahlung.

\section{Temperatur und Entropie}
Die Hawking-Temperatur eines Schwarzschild-Lochs ist
\begin{equation}
T_H=\frac{\hbar c^3}{8\pi G M k_B}.
\end{equation}
Die Bekenstein-Hawking-Entropie lautet
\begin{equation}
S_{BH}=\frac{k_B c^3 A}{4G\hbar}.
\end{equation}
Beide Formeln enthalten zugleich $G$, $c$ und $\hbar$. Sie liegen also genau im Überschneidungsbereich von Gravitation, Relativität und Quantentheorie.

\section{Bedeutung}
Schwarze Löcher zeigen, dass Raumzeitgeometrie thermodynamische und quantenmechanische Eigenschaften besitzt. Auch dies wählt keine fertige Theorie aus, verstärkt aber die Notwendigkeit eines quantengravitativen Abschlusses.

\chapter{Der geschlossene Hauptbeweis}
\section{Annahmen}
Wir fassen die Voraussetzungen zusammen.

\begin{axiom}[Relativistische Quantenmaterie]
Materie wird durch lokale relativistische Quantenfelder beschrieben. Energie und Impuls sind durch einen Energie-Impuls-Operator gegeben.
\end{axiom}

\begin{axiom}[Dynamische Raumzeit]
Gravitation wird durch eine dynamische Lorentz-Metrik beschrieben. Die Metrik besitzt propagierende Freiheitsgrade.
\end{axiom}

\begin{axiom}[Universelle Kopplung]
Gravitation koppelt universell an Energie und Impuls.
\end{axiom}

\begin{axiom}[Unitäre gemeinsame Entwicklung]
Der vollständige Theorieabschluss darf keine unkontrollierte Verletzung der quantenmechanischen Zustandsentwicklung einführen.
\end{axiom}

\section{Beweisschritte}
\begin{lemma}[Quantenmaterie erzeugt operatorwertige Quellen]
Unter Annahme 1 ist der Energie-Impuls-Tensor der Materie im Allgemeinen operatorwertig und besitzt Fluktuationen.
\end{lemma}
\begin{proof}
In der Feldquantisierung werden Felder und ihre kanonischen Impulse zu Operatoren. Da $T_{\mu\nu}$ aus diesen Feldern gebildet wird, wird auch $T_{\mu\nu}$ operatorwertig. Erwartungswerte bestimmen den Zustand nicht vollständig.
\end{proof}

\begin{lemma}[Gravitation besitzt eigene Freiheitsgrade]
Unter Annahme 2 enthält die linearisierte Gravitation propagierende Freiheitsgrade.
\end{lemma}
\begin{proof}
Die Entwicklung $g_{\mu\nu}=\eta_{\mu\nu}+\kappa h_{\mu\nu}$ führt nach Eichfixierung auf eine Wellengleichung für die physikalischen Komponenten. Diese Komponenten beschreiben Gravitationswellen.
\end{proof}

\begin{lemma}[Universelle Kopplung macht die Wechselwirkung wechselseitig]
Unter Annahme 3 beeinflusst Quantenmaterie die Geometrie, und dynamische Geometrie beeinflusst Quantenmaterie.
\end{lemma}
\begin{proof}
Die Materiewirkung hängt von der Metrik ab, daher beeinflusst Geometrie die Materiedynamik. Zugleich ist die Variation nach der Metrik durch $T_{\mu\nu}$ bestimmt, daher beeinflusst Materie die Geometrie.
\end{proof}

\begin{lemma}[Eine vollständig klassische Geometrie ist nur Näherung]
Eine Geometrie, die nur klassisch bleibt und nur an Erwartungswerte quantisierter Materie koppelt, enthält im Allgemeinen nicht die vollständige Quanteninformation der Quelle.
\end{lemma}
\begin{proof}
Die Abbildung eines Quantenzustands auf Erwartungswerte von $\hat T_{\mu\nu}$ ist nicht injektiv. Unterschiedliche Zustände können dieselbe mittlere Quelle besitzen. Eine klassische Geometrie, die nur von dieser mittleren Quelle abhängt, kann die verschiedenen Zustände nicht vollständig unterscheiden.
\end{proof}

\section{Hauptsatz}
\begin{satz}[Notwendigkeit der Quantengravitation]
Unter den vier Annahmen ist ein gemeinsamer, geschlossener Theorieabschluss nur möglich, wenn die gravitativen Freiheitsgrade quantenfähig sind. Eine vollständig klassische Gravitation ist höchstens eine Näherung, nicht die Endform.
\end{satz}
\begin{proof}
Nach Annahme 1 besitzt Materie quantenmechanische Zustandsstruktur, und ihre Energie-Impuls-Größen sind im Allgemeinen operatorwertig. Nach Annahme 2 besitzt Gravitation eigene dynamische Freiheitsgrade. Nach Annahme 3 koppeln diese Freiheitsgrade universell an Energie und Impuls und wirken zugleich auf die Materiedynamik zurück.

Würde Gravitation vollständig klassisch bleiben, könnte sie höchstens an klassische Daten wie Erwartungswerte koppeln. Nach dem vorherigen Lemma reichen solche Daten nicht aus, um die vollständige Quantenstruktur der Quelle zu enthalten. Außerdem können propagierende gravitative Freiheitsgrade mit Quantenmaterie Energie und Impuls austauschen. Ein solcher Austausch ist im vollständigen Rahmen eine gemeinsame Zustandsentwicklung. Bleiben die gravitativen Freiheitsgrade außerhalb der Quantenzustandsbeschreibung, entsteht keine geschlossene gemeinsame Theorie, sondern nur eine semiklassische Näherung.

Damit verlangt der gemeinsame Abschluss, dass auch die gravitativen Freiheitsgrade quantenfähig sind. Genau diese Forderung ist die Notwendigkeit der Quantengravitation.
\end{proof}

\begin{korollar}[Klassische Grenze]
Die allgemeine Relativitätstheorie bleibt als klassische Grenztheorie erhalten, wenn quantengravitative Korrekturen vernachlässigbar sind.
\end{korollar}
\begin{proof}
Im Grenzfall großer Wirkung gegenüber $\hbar$ oder schwacher Quantenfluktuationen führen stationäre Beiträge wieder auf die klassischen Feldgleichungen. Daher wird die klassische Theorie nicht verworfen, sondern als Grenzfall eingeordnet.
\end{proof}

\chapter{Was abgeschlossen ist und was offen bleibt}
\section{Abgeschlossen}
Abgeschlossen ist der Konsistenzbeweis: Quantenmaterie, dynamische Raumzeit, universelle Kopplung und unitäre gemeinsame Entwicklung verlangen quantenfähige gravitative Freiheitsgrade. Damit ist Quantengravitation als notwendiger Abschluss der bekannten Prinzipien hergeleitet.

\section{Offen}
Offen bleibt die Auswahl der vollständigen Hochenergie-Theorie. Die bekannten Forschungsrichtungen geben unterschiedliche Antworten auf die konkrete mathematische Ausarbeitung. Dieses Skript wählt keine davon endgültig aus.

Offen bleiben insbesondere:
\begin{enumerate}
\item die vollständige nichtstörungstheoretische Definition,
\item die Behandlung singulärer Raumzeiten,
\item die mikroskopische Deutung der Schwarze-Loch-Entropie,
\item die vollständige experimentelle Prüfung im Planck-Bereich,
\item die Verbindung zur beobachteten Kosmologie im frühen Universum.
\end{enumerate}

\section{Warum der Beweis dennoch rund ist}
Der Beweis ist rund, weil er genau das zeigt, was aus bekannter Mathematik und bekannter Physik zwingend folgt: Nicht die fertige Auswahl einer bestimmten Hochenergie-Theorie, sondern die Notwendigkeit quantenfähiger Gravitation. Er ist abgeschlossen, weil jede verwendete Voraussetzung genannt und jeder Übergang begründet wurde.



\chapter{Newtonsche Grenze und Gravitationswellen}
\section{Schwaches Feld}
Im schwachen Feld schreibt man
\begin{equation}
g_{00}=-(1+2\Phi/c^2),\qquad g_{ij}=(1-2\Phi/c^2)\delta_{ij}.
\end{equation}
Für langsam bewegte Materie ist $T_{00}\approx \rho c^2$. Die $00$-Komponente der Einstein-Gleichung reduziert sich in führender Ordnung auf die Poisson-Gleichung
\begin{equation}
\Delta \Phi=4\pi G\rho.
\end{equation}
Damit enthält die allgemeine Relativitätstheorie die Newtonsche Gravitation als Grenzfall.

\begin{lemma}[Korrespondenzprinzip]
Jede tragfähige Quantengravitation muss bei schwachen Feldern, niedrigen Energien und großen Wirkungen gegenüber $\hbar$ sowohl die Newtonsche Gravitation als auch die allgemeine Relativitätstheorie als Grenzfälle enthalten.
\end{lemma}
\begin{proof}
Die empirisch bestätigten Grenzfälle dürfen im gemeinsamen Abschluss nicht verschwinden. Da die Einstein-Gleichung im schwachen statischen Grenzfall die Poisson-Gleichung liefert, muss jede tiefere Theorie diese Näherung reproduzieren.
\end{proof}

\section{Gravitationswellen als dynamische Freiheitsgrade}
Gravitationswellen sind keine Rechenartefakte. Sie sind Lösungen der Feldgleichungen und wurden beobachtet. Ihre Existenz bedeutet, dass die Geometrie selbst Energie übertragen kann. Ein Objekt, das Energie übertragen und mit Quantenmaterie wechselwirken kann, ist in einem vollständigen Quantenrahmen kein bloßer äußerer Parameter.

\begin{merkkasten}
\textbf{Schlüsselpunkt.} Die Gravitation besitzt eigene bewegliche Freiheitsgrade. Genau diese Freiheitsgrade sind es, deren Quantenfähigkeit im Hauptsatz folgt.
\end{merkkasten}

\chapter{Erhaltungsgrößen und Symmetrien}
\section{Noether-Gedanke}
Symmetrien und Erhaltungsgrößen sind in der Physik untrennbar verbunden. Zeittranslationssymmetrie führt zu Energieerhaltung, Raumtranslationssymmetrie zu Impulserhaltung, Rotationssymmetrie zu Drehimpulserhaltung.

In gekrümmter Raumzeit wird diese Ordnung lokal und kovariant. Die Bedingung
\begin{equation}
\nabla_\mu T^{\mu\nu}=0
\end{equation}
ist Ausdruck lokaler Energie-Impuls-Erhaltung im geometrischen Rahmen.

\section{Diffeomorphismeninvarianz}
Die allgemeine Relativitätstheorie ist invariant unter glatten Koordinatentransformationen. Diese Invarianz bedeutet nicht, dass alles beliebig ist. Sie bedeutet, dass physikalische Aussagen unabhängig von der Wahl der Koordinaten formuliert werden müssen.

Diese Symmetrie ist tief. Sie ist zugleich Grund dafür, dass die kanonische Theorie Zwangsbedingungen enthält. Die Zeitentwicklung ist nicht einfach Bewegung in einem äußeren Zeitparameter, sondern Teil der dynamischen Geometrie.

\begin{satz}[Symmetrie als Konsistenzbedingung]
Die Diffeomorphismeninvarianz zwingt die Gravitationsgleichungen, Quellen nur in kovariant erhaltener Form aufzunehmen.
\end{satz}
\begin{proof}
Die kontrahierte Bianchi-Identität ist eine geometrische Folge der Diffeomorphismeninvarianz. Da die linke Seite der Feldgleichung kovariant divergenzfrei ist, muss auch die Quelle kovariant divergenzfrei sein.
\end{proof}

\chapter{Quantisierung freier Felder genauer betrachtet}
\section{Vom harmonischen Oszillator zum Feld}
Ein freies Feld kann als unendliche Menge gekoppelter Schwingungsmoden verstanden werden. Jede Mode verhält sich wie ein harmonischer Oszillator. Die Quantisierung des harmonischen Oszillators führt auf Erzeugungs- und Vernichtungsoperatoren. Für ein Feld entstehen daraus Teilchenanregungen.

Für ein skalares Feld gilt schematisch
\begin{equation}
\phi(x)=\int \frac{\diff^3p}{(2\pi)^3}\frac{1}{\sqrt{2E_p}}\left(a_p e^{-ipx/\hbar}+a_p^\dagger e^{ipx/\hbar}\right).
\end{equation}
Die Operatoren erfüllen Vertauschungsrelationen. Dadurch wird aus einer klassischen Feldamplitude ein quantisiertes Feld.

\section{Übertragung auf die Gravitation}
Im linearisierten Grenzfall besitzt auch die Metrikstörung wellenartige Moden. Ihre Quantisierung führt zu Quanten der gravitativen Welle. Der Name dieser Quanten ist nicht für den Beweis entscheidend. Entscheidend ist die Struktur: Eine dynamische Welle mit Energie, Impuls und Wechselwirkung wird in der Quantentheorie durch Feldoperatoren beschrieben.

\begin{lemma}[Feldmodus-Argument]
Jede propagierende lineare Feldmode, die mit Quantenmaterie wechselwirkt und Energie austauscht, verlangt im vollständigen Quantenrahmen eine operatorielle Beschreibung.
\end{lemma}
\begin{proof}
Eine propagierende Mode kann angeregt, abgeschwächt und mit anderen Quantensystemen gekoppelt werden. In der Quantentheorie werden solche Vorgänge durch Übergangsamplituden zwischen Zuständen beschrieben. Eine rein klassische Beschreibung kann diese Übergänge höchstens als Näherung darstellen.
\end{proof}

\chapter{Messung, Rückwirkung und Verschränkung}
\section{Rückwirkung}
Messung in der Quantentheorie ist nicht bloß Ablesen. Eine Wechselwirkung zwischen Messgerät und System kann den Zustand verändern. Wenn das Messgerät über Masse und Energie verfügt, koppelt es auch gravitativ.

Eine vollständig gemeinsame Theorie muss daher nicht nur die Wirkung der Geometrie auf Materie, sondern auch die Rückwirkung der Materie auf die Geometrie behandeln.

\section{Verschränkung als Prüfstein}
Wenn zwei Quantensysteme über ein vermittelndes Feld wechselwirken, kann dieses Feld Korrelationen erzeugen. Wird durch gravitative Wechselwirkung Verschränkung zwischen Materiesystemen erzeugt, dann ist dies ein starkes Indiz dafür, dass die vermittelnde Struktur selbst nicht vollständig klassisch sein kann.

Für den allgemeinen Beweis genügt die schwächere Aussage: Ein Feld, das im vollständigen Rahmen Quanteninformation zwischen Quantensystemen vermittelt, darf nicht nur klassischer Buchhalter mittlerer Werte sein.

\begin{satz}[Vermittlungsprinzip]
Ein vollständig klassischer Vermittler kann im allgemeinen Fall keine vollständige quantenmechanische Zustandsentwicklung zweier wechselwirkender Quantensysteme ersetzen.
\end{satz}
\begin{proof}
Eine quantenmechanische Wechselwirkung wird durch gemeinsame Zustände und unitäre Entwicklung beschrieben. Klassische Daten erfassen nur einen Teil der quantenmechanischen Struktur. Sie enthalten im Allgemeinen weder Phasenbeziehungen noch alle Korrelationen. Daher kann ein rein klassischer Vermittler nur eine Näherung bilden.
\end{proof}

\chapter{Renormierung und ihre Bedeutung}
\section{Warum Unendlichkeiten auftreten}
Quantenfeldtheorien erzeugen in Störungsrechnungen häufig divergente Ausdrücke. Renormierung ordnet diese Ausdrücke so, dass endliche, messbare Größen entstehen. In renormierbaren Theorien genügt eine endliche Zahl von Parametern zur Festlegung der Theorie.

Bei der Gravitation treten wegen der Dimension der Newton-Konstante bei hohen Energien unendlich viele Gegenstücke höherer Krümmungsterme auf. Deshalb ist die störungstheoretische Quantisierung der allgemeinen Relativitätstheorie nicht als einfache renormierbare Hochenergie-Theorie abgeschlossen.

\section{Warum dies den Hauptsatz nicht widerlegt}
Die fehlende gewöhnliche Renormierbarkeit bedeutet nicht, dass gravitative Freiheitsgrade nicht quantenfähig sein müssen. Sie bedeutet nur, dass die naive störungstheoretische Fortsetzung nicht die endgültige Hochenergieform liefert.

Als wirksame Feldtheorie bleibt die Quantisierung bei niedrigen Energien sinnvoll. Für die höchste Energieordnung wird eine tiefere mathematische Struktur benötigt.

\begin{merkkasten}
\textbf{Wichtige Trennung.} Die Notwendigkeit der Quantengravitation ist stärker als die Auswahl einer bestimmten Quantengravitationstheorie. Der Hauptsatz beweist die Notwendigkeit des quantenfähigen Abschlusses, nicht die endgültige Wahl seines Formalismus.
\end{merkkasten}

\chapter{Kausalstruktur und Lichtkegel}
\section{Lichtkegel}
Die Lorentz-Metrik bestimmt an jedem Ereignis einen Lichtkegel. Er trennt mögliche kausale Wirkungen von raumartigen Trennungen. Diese Struktur ist nicht nachträglich aufgesetzt, sondern Bestandteil der Geometrie.

Wenn die Geometrie dynamisch ist, ist auch die Kausalstruktur dynamisch. In einer Theorie, in der Geometrie quantenfähig ist, wird damit auch die Beschreibung kausaler Struktur subtiler.

\section{Grenzforderung}
Trotzdem muss die klassische Kausalstruktur in geeigneten Grenzfällen zurückkehren. Die Theorie darf nicht willkürlich werden. Sie muss die bekannten Lichtkegelstrukturen dort reproduzieren, wo klassische Raumzeit eine gute Näherung ist.

\begin{lemma}[Kausalgrenze]
Ein quantengravitativer Abschluss muss im klassischen Grenzfall die kausale Ordnung der Lorentz-Geometrie zurückgeben.
\end{lemma}
\begin{proof}
Die Lorentz-Geometrie ist empirisch bestätigt und folgt aus der klassischen Grenze. Jede tiefere Theorie, die diese Grenze nicht reproduziert, würde die bekannten Niederenergiephänomene verfehlen.
\end{proof}

\chapter{Kosmologische Anwendung}
\section{Frühes Universum}
Im sehr frühen Universum waren Energiedichten hoch und Längenskalen klein. Daher ist zu erwarten, dass Quantenfelder und Gravitation dort nicht getrennt behandelt werden können.

Kosmologische Anfangszustände, Vakuumfluktuationen und Strukturbildung verweisen auf eine gemeinsame Behandlung von Quantenphysik und dynamischer Raumzeit.

\section{Klassische Gleichungen als Näherung}
Die Friedmann-Gleichungen beschreiben die großskalige kosmische Entwicklung. Sie folgen aus der allgemeinen Relativitätstheorie unter Symmetrieannahmen. Im frühen Grenzbereich müssen jedoch Quantenkorrekturen erwartet werden.

\begin{equation}
H^2=\frac{8\pi G}{3}\rho-\frac{kc^2}{a^2}+\frac{\Lambda c^2}{3}.
\end{equation}
Diese Gleichung bleibt wichtig, aber sie ist nicht der letzte Satz, wenn $\rho$ extrem groß und quantenmechanische Fluktuationen wesentlich werden.

\chapter{Vergleich möglicher Ausarbeitungen}
\section{Gemeinsamer Kern}
Verschiedene Ansätze zur Quantengravitation unterscheiden sich in ihrer mathematischen Ausgestaltung. Ihr gemeinsamer Kern ist jedoch die Einsicht, dass klassische Geometrie und Quantentheorie nicht einfach unvermittelt nebeneinander stehen können.

Einige Ansätze beginnen bei der Quantisierung geometrischer Größen. Andere beginnen bei erweiterten Objekten, kausalen Mengen, thermodynamischen Argumenten oder Informationsprinzipien. Für den Hauptsatz dieses Skripts ist diese Vielfalt kein Problem. Der Hauptsatz liegt eine Ebene darunter: Er begründet, warum ein quantenfähiger Abschluss überhaupt erforderlich ist.

\section{Auswahlproblem}
Welche vollständige Ausarbeitung zutrifft, muss durch mathematische Konsistenz, Grenzfallprüfung und empirische Unterscheidbarkeit entschieden werden. Das Auswahlproblem bleibt offen, aber die Notwendigkeit des quantenfähigen Abschlusses bleibt bestehen.

\chapter{Schließung der Beweiskette}
\section{Keine Lücke zwischen Quelle und Geometrie}
Der wesentliche Punkt ist die Quelle. Die Geometrie koppelt an Energie und Impuls. In der Quantenfeldtheorie sind Energie und Impuls Bestandteile der operatoriellen Zustandsstruktur. Eine Theorie, die diese Quelle vollständig ernst nimmt, kann die Geometrie nicht dauerhaft außerhalb des gemeinsamen Zustandsrahmens halten.

\section{Keine Lücke zwischen Welle und Quant}
Der zweite Punkt ist die Welle. Die Gravitation besitzt propagierende Wellenfreiheitsgrade. In einer Quantentheorie werden propagierende Wellen mit Energieaustausch durch quantisierte Feldmoden beschrieben. Auch hier ist eine rein klassische Endform nur Näherung.

\section{Keine Lücke zwischen Grenze und Abschluss}
Der dritte Punkt ist der Grenzfall. Die klassische allgemeine Relativitätstheorie muss erhalten bleiben, aber als Grenzfall. Genau dies leistet ein quantenfähiger Abschluss: Er enthält die klassische Theorie dort, wo Quanteneffekte vernachlässigbar sind, und überschreitet sie dort, wo Materie, Raumzeit und $\hbar$ gemeinsam wesentlich werden.

\begin{satz}[Abschlusssatz]
Die drei Forderungen -- vollständige Quellenbeschreibung, dynamische gravitative Wellenfreiheitsgrade und Erhaltung des klassischen Grenzfalls -- schließen zusammen die rein klassische Gravitation als Endform aus und begründen die Notwendigkeit der Quantengravitation.
\end{satz}
\begin{proof}
Die vollständige Quellenbeschreibung verlangt die Berücksichtigung der Quantenstruktur von Energie und Impuls. Die dynamischen gravitativen Wellenfreiheitsgrade verlangen eine Zustandsbeschreibung der Geometrie im gemeinsamen Austausch mit Materie. Die klassische Grenze verlangt, dass die bekannte allgemeine Relativitätstheorie als Näherung erhalten bleibt. Eine rein klassische Gravitation erfüllt höchstens die dritte Forderung, aber nicht die ersten beiden im vollständigen Sinn. Daher ist sie nicht Endform. Ein quantenfähiger gravitativer Abschluss erfüllt alle drei Forderungen dem Prinzip nach. Damit ist die Notwendigkeit der Quantengravitation bewiesen.
\end{proof}

\chapter{Übungen}
\section*{Aufgaben}
\addcontentsline{toc}{section}{Aufgaben}
\begin{enumerate}
\item Zeigen Sie aus der kontrahierten Bianchi-Identität, dass die Einstein-Gleichung die kovariante Erhaltung des Energie-Impuls-Tensors verlangt.
\item Erklären Sie, warum ein Erwartungswert einen Quantenzustand im Allgemeinen nicht eindeutig bestimmt.
\item Leiten Sie die Eichtransformation der linearisierten Metrikstörung aus einer infinitesimalen Koordinatentransformation her.
\item Begründen Sie, warum ein masseloses Feld mit Drehimpulszahl zwei an einen erhaltenen Energie-Impuls-Tensor koppeln muss.
\item Erklären Sie, warum semiklassische Gravitation als Näherung sinnvoll, aber nicht als vollständiger Abschluss ausreichend ist.
\item Geben Sie die Planck-Länge, Planck-Zeit und Planck-Masse an und erklären Sie ihre Bedeutung.
\item Beschreiben Sie, warum schwarze Löcher auf einen Zusammenhang von Gravitation, Quantentheorie und Thermodynamik hinweisen.
\item Formulieren Sie den Hauptsatz dieses Skripts in eigenen Worten.
\end{enumerate}

\section*{Musterlösungen in Kurzform}
\addcontentsline{toc}{section}{Musterlösungen in Kurzform}
\begin{enumerate}
\item Anwendung von $\nabla^\mu G_{\mu\nu}=0$ und $\nabla^\mu g_{\mu\nu}=0$ auf die Einstein-Gleichung liefert $\nabla^\mu T_{\mu\nu}=0$.
\item Ein Erwartungswert ist nur eine Zahl. Der Zustand enthält zusätzlich Varianzen, Phasenbeziehungen und Korrelationen.
\item Die Koordinatenänderung verschiebt die Metrikstörung um symmetrisierte Ableitungen des Verschiebungsfeldes.
\item Eichinvarianz der Kopplung verlangt Erhaltung der Quelle; diese Quelle ist der Energie-Impuls-Tensor.
\item Sie verwendet Erwartungswerte und verfehlt damit im Allgemeinen Fluktuationen und vollständige Zustandsinformation.
\item $\ell_P=\sqrt{\hbar G/c^3}$, $t_P=\sqrt{\hbar G/c^5}$, $m_P=\sqrt{\hbar c/G}$. Sie markieren die gemeinsame Grenzskala von Gravitation und Quantentheorie.
\item Temperatur und Entropie schwarzer Löcher enthalten $G$, $c$ und $\hbar$ zugleich.
\item Wenn Quantenmaterie und dynamische Raumzeit vollständig gemeinsam beschrieben werden sollen, müssen auch gravitative Freiheitsgrade quantenfähig sein.
\end{enumerate}

\chapter*{Schlussformel}
\addcontentsline{toc}{chapter}{Schlussformel}
Die Beweiskette des dritten Semesters lautet:

\begin{merkkasten}
\raggedright
Quantenmaterie besitzt operatorwertige Energie und Impuls.

Energie und Impuls krümmen dynamische Raumzeit.

Dynamische Raumzeit besitzt eigene propagierende Freiheitsgrade.

Diese Freiheitsgrade koppeln universell und wechselseitig an Quantenmaterie.

Eine rein klassische Geometrie kann die vollständige Quantenstruktur dieser Kopplung nicht abschließend tragen.

Also müssen gravitative Freiheitsgrade im gemeinsamen Abschluss quantenfähig sein.

Das ist die Notwendigkeit der Quantengravitation.
\end{merkkasten}

\vspace{1em}
\begin{center}
\textbf{Quantengravitation ist der notwendige Konsistenzabschluss von Quantenmaterie und dynamischer Raumzeit.}
\end{center}

q.e.d.\par
Ingolf Lohmann

\end{document}
